\section{Teoretické východiská}
\label{sec:teoria}

% Kapitola vychádza z odbornej literatúry (zdroje sú v references.bib).
% Rozsah má byť najviac 3 strany. Pred odovzdaním skontrolujte, že každé
% tvrdenie zodpovedá obsahu citovaného zdroja.

\subsection{Behaviorálna biometria}

Biometrická autentifikácia overuje alebo určuje totožnosť človeka podľa jeho
telesných alebo behaviorálnych vlastností. Behaviorálna biometria vychádza zo
spôsobu, akým človek koná, a dynamika písania na klávesnici patrí k jej
najľahšie dostupným podobám: nepotrebuje špeciálny snímač, stačí bežná
klávesnica a údaje sa dajú zbierať počas písania~\cite{roy2022,shadman2023}.
Prehľad aktuálnej literatúry poskytuje systematický prehľad zameraný na
modely dynamiky písania z rokov 2017 až 2022~\cite{roy2022}, ktorý porovnáva
použité metódy, nastavenia experimentov aj otvorené problémy, a prehľadová
štúdia o pojmoch, technikách a aplikáciách, ktorá sa venuje aj dátovým súborom,
písaniu na dotykových zariadeniach a použitiu mimo autentifikácie~\cite{shadman2023}.

Rozlišujú sa dve úlohy: \emph{verifikácia} (overenie, že človek je tým, za koho
sa vydáva) a \emph{identifikácia} (určenie, ktorému zo známych ľudí sa
záznam najviac podobá). Táto práca rieši identifikáciu. Podľa typu písaného
textu sa ďalej rozlišuje písanie \emph{pevného} textu (všetci píšu to isté,
ako v tejto práci) a písanie \emph{voľného} textu, ktoré umožňuje priebežné
overovanie počas bežnej práce na počítači~\cite{shadman2023}.

\subsection{Príznaky dynamiky písania}

Základom sú časy stlačenia a uvoľnenia klávesov. Z nich sa odvádzajú dva
druhy príznakov: doba držania klávesy (čas medzi stlačením a uvoľnením)
a intervaly medzi po sebe idúcimi klávesmi, ktoré sa v literatúre nazývajú aj
letové časy (\emph{flight time})~\cite{kaluarachchi2023}. Autori práce o súbore
DEFT ukazujú, že doplnenie letových časov o vzdialenosť medzi klávesmi
vytvára príznakovú množinu nezávislú od zariadenia,
s presnosťou nad 99\,\% a chybovosťou EER pod 10\,\% na troch zariadeniach~\cite{kaluarachchi2023}.
Zaujímavým zistením je, že aj príznaky založené len na časoch, bez znalosti
obsahu textu, si zachovávajú uspokojivý výkon pri verifikácii~\cite{stragapede2023}.
Práve takéto časové príznaky (doby držania a intervaly) používa aj riešenie
v tejto práci (podkapitola~\ref{sec:priznaky}).

\subsection{Metódy vyhodnotenia}

Pri klasifikácii sa používajú jednoduché štatistické miery vzdialenosti aj
metódy strojového učenia. Porovnanie klasických algoritmov strojového učenia
a techník výberu príznakov (filtrovacích a obalových) na úlohe autentifikácie
podľa dynamiky písania uvádza štúdia~\cite{teotia2023}. Pri voľnom texte sa
úspešne používajú aj hlboké neurónové siete: rekurentná sieť TypeNet
dosiahla chybovosť EER 2,2\,\% na fyzickej klávesnici a 9,2\,\% na dotykovej
a pri nárasti počtu používateľov až do 100\,000 sa chyba zvýšila len mierne~\cite{acien2021},
a kombinácia konvolučnej a rekurentnej siete (CNN a GRU) dosiahla na voľnom
texte výsledky porovnateľné s najlepšími predchádzajúcimi prístupmi~\cite{chang2021}.
Autori porovnania z roku 2026 zdôrazňujú, že výsledky veľmi závisia od výberu
dátového súboru a spôsobu vyhodnotenia, že hlboké učenie vedie pri voľnom texte
a že tvrdenia o kvantových prístupoch zatiaľ nie sú podložené~\cite{gundogan2026}.
Aj preto vznikla verejná výzva KVC, ktorá porovnáva systémy na rovnakých dátach
a okrem presnosti hodnotí aj spravodlivosť voči demografickým skupinám~\cite{stragapede2023}.

Verifikačné systémy sa najčastejšie hodnotia mierou EER (\emph{equal error
rate}), teda chybovosťou, pri ktorej je podiel nesprávne prijatých pokusov
rovnaký ako podiel nesprávne odmietnutých~\cite{acien2021,stragapede2023}.
Pri identifikácii na malom počte ľudí je
prirodzenejšie sledovať, či je správny človek na prvom mieste (top-1) alebo
medzi prvými $k$ (top-$k$); tieto metriky používa aj táto práca
(podkapitola~\ref{sec:protokol}). Nie sú priamo porovnateľné s EER, preto sa
výsledky tejto práce nedajú jednoducho porovnávať s číslami z citovaných
štúdií, ktoré používajú iné dáta a iný protokol.

\subsection{Súkromie a ochrana osobných údajov}

Dynamika písania je citlivý údaj. Behaviorálne príznaky môžu prezradiť aj
pohlavie, vek, dominantnú ruku či emocionálny stav písajúceho~\cite{baig2022},
a preto sa navrhujú protokoly, v ktorých server nikdy nevidí biometrické šablóny
v otvorenej podobe, napríklad s využitím homomorfného šifrovania~\cite{baig2022}.
Pri právnom posúdení treba brať do úvahy, že terminológia GDPR rozlišuje
rozpoznávanie osoby a len klasifikáciu jej vlastností (napríklad tzv. mäkkých
znakov) a že druhý prípad je chránený menej, čo vytvára riziká pre súkromie~\cite{bisztray2021}. Tieto
úvahy sú dôvodom, prečo sa pri zbere údajov v tejto práci uplatňuje informovaný
súhlas, minimalizácia údajov a lokálne úložisko (podkapitola~\ref{sec:etika}).
